\chapter*{2025版序}
\addcontentsline{toc}{chapter}{2025版序}
\SourceParagraph{首先欢迎大家来到南赫学院！！！}
\SourceParagraph{本手册是由一群现在大一、大二的同学自发建设的项目，旨在分享各类南赫学院相关的信息，帮助新生适应大学生活，进行生涯规划。经过前期筹划，生存手册于 2025 年 6 月 21 日开始编写，8 月 20 号完成初稿的建设。手册的内容以参与编写的成员的分享为主，辅以少量约稿，主要包含培养体系、学习建议以及生涯规划这三大板块。手册草草成稿，缺漏谬误之处在所难免，尽管部分错误已经被修正，目前的版本尚留有许多需要完善、修订的地方，故将其取名为“难喝生存手册 v0”，以期未来增补。目前成型的章节汇聚了学长学姐们的一些经验教训，或许能对大家有所裨益。}
\SourceParagraph{2023 年 9 月，我以第一届本科生的身份来到了刚刚建成的南赫学院。初入大学，我曾就大气专业就业升学相关的问题感到焦虑惶恐，质疑自己在专业院校的选择上受到了世俗价值观的裹挟，又忧虑于埋头学习艰涩繁难的知识会限制自己未来职业发展的道路，于是时常求教身边的良师益友，与许多优秀真诚的学长学姐们相识。他们给予我的支持、建议与期待，帮助我走出困顿与迷茫，并渐渐认识到，不完美是生活的常态，不必过于焦虑自身的天赋和未来的不确定性。我希望以这份手册为起点，搭建一个平台，吸引更多人分享经验与感悟，帮助、鼓励并陪伴未来的同学共同面对挑战。我个人能力有限、经验有限，只有更多人参与，信息才能更全面、贴近大家的需要。}
\SourceParagraph{必须承认，南赫学院尚处起步阶段，确实存在各种问题。不过我们不会在此抨击现行体制或政策。无论未来它会发展成怎样，我们都希望它变得更好。确定编写手册这一想法后，我与许多人进行了交流，得到了广泛支持，甚至有伙伴提前行动。正是这些志同道合的伙伴，让手册从一个模糊的设想到逐步落地为今日的文字。在未来，我们期待你成为手册的作者，用你的亲身经历去丰富、修订和传承这份手册。}
\SourceBlankLine[2]
\begin{flushright}
\begin{tabular}{@{}c@{}}
2025 年 8 月 20 日\\[1em]
编者按
\end{tabular}
\end{flushright}


\addcontentsline{toc}{chapter}{2026版序}
\chapter*{2026版序}
\SourceParagraph{首先，欢迎大家来到南赫学院。}

\SourceParagraph{距离《难喝生存手册 v0》第一次与大家见面，已经过去了一年。去年，我们怀着一些朴素的想法开始编写这本手册：希望把已经知道的信息整理下来，把走过的弯路、踩过的坑和偶尔获得的一点经验分享给后来的人，也希望在这个仍然年轻的学院里，留下些什么。如今再回头看，这本手册当然远称不上完善，但它确实从最初几个人之间的设想，慢慢变成了一份可以被阅读、讨论、修改和继续传递下去的文字。更令人欣喜的是，在新一届编写者的加入下，越来越多同学愿意参与其中，用自己的经历、思考和感悟丰富这份手册。它也因此不再只是少数人的记录，而逐渐成为属于南赫学生共同建设的一份资料。于是，我们决定继续把它写下去。}

\SourceParagraph{我们之所以把它叫作“生存手册”，其实多少带有一点玩笑意味，但又并不完全是玩笑。}

\SourceParagraph{刚来到这里的时候，很多人或许都曾有过失望、不适应甚至抱怨。大学和想象中的样子并不总是一致，学院本身也仍处在建设和调整之中。课程安排、培养体系、校园生活、地理位置，乃至一些看起来微不足道的小事，都可能在某个时刻让人感到烦躁。我们当然可以期待身边的环境变得更好，也应该在有能力的时候提出意见、推动改变，但很多时候，我们确实无法立即改变周围的一切，更无法要求所有的人和事都按照自己的期待运行。}

\SourceParagraph{所以，“生存”在这里并不是消极地忍耐，更不是告诉大家“既然来了就只能接受”。它更接近一种主动适应环境的能力：在有限的条件里获取尽可能充分的信息，认识自己真正想要什么，学会做选择、解决问题，也尽量让自己的大学生活过得更从容一些。我们未必能够决定身边的一切，却可以努力决定自己怎样回应这些事情。所谓“既来之，则安之”，或许不应该被理解成认命，而应该是承认现实之后，仍然认真地生活，并努力把自己能够做好的事情做好。}

\SourceParagraph{这也是我们继续编写这本手册最重要的原因之一。}

\SourceParagraph{最开始产生这个想法，是因为我们看到其他学院已经有了各种各样的经验分享、课程指南和生存手册。而南赫学院成立时间并不长，许多经验尚未形成稳定的传承，很多信息也只能依靠同学之间口耳相传。于是我们想，既然总要有人去经历这些事情，不如把其中有价值的部分留下来。也许若干年以后，课程会改变，培养方案会调整，学院会变得和现在很不一样，但至少这些文字可以记录下我们这一代学生怎样学习、怎样迷茫、怎样试错，又怎样慢慢找到自己的方向。}

\SourceParagraph{如果这本手册能够让后来的人少花一点时间寻找信息，少踩几个前人已经踩过的坑；能够让某个正因为选课、专业、升学或未来而焦虑的人发现，“原来以前也有人遇到过一样的问题”；或者哪怕只是让大家读到某一段文字时有所感触，那么它存在的意义就已经足够了。}

\SourceParagraph{我们也始终希望，这不是一本由少数人定义“大学应该怎样度过”的手册。每个人的经历都非常有限，同一件事情，从不同的人那里往往会得到完全不同的答案。因此，比起给出一个所谓的“标准答案”，我们更希望它成为一个交流和分享的平台。有人分享课程经验，有人谈科研和竞赛，有人讨论保研、留学和就业，也有人记录自己对于校园生活的感受。观点可以不同，选择当然也可以不同。只有越来越多的人愿意把自己的经历留下来，这本手册才可能真正接近大家所需要的样子。}

\SourceParagraph{当然，我们也知道，有些同学来到这里以后，可能会因为种种问题怀疑自己当初的选择，甚至想象：如果当时去了另一所学校、另一个学院、另一个专业，大学生活是不是会完全不同？}

\SourceParagraph{这种想法非常正常。但后来我越来越觉得，我们很难比较自己真正走过的路与一条从未走过的路，也很容易在失望的时候，不自觉地美化那个没有选择的可能。很多焦虑、竞争、信息差和对于未来的不确定，并不只属于南赫，也并不只属于某一个学院。换一条路当然可能遇见不同的风景，但也同样会遇见属于那条路的问题。我们没有必要因为这里存在不足就否认自己的感受，也没有必要因此认定另一种选择必然更加完美。}

\SourceParagraph{至少在我自己的经历里，我很幸运地遇到了许多值得尊敬和感谢的人。这学院里的领导和老师，他们中的很多人都真诚地关心学生，也在各自能够做到的范围里尽力解决问题、改善环境、推动学院向前发展。一个刚刚建立不久的学院必然需要时间形成成熟的培养体系、传统和文化，而很多改变本身也不可能在一夜之间发生。}

\SourceParagraph{因此，我们当然可以指出问题，也应该提出合理的批评和建议；但在看到问题之外，也许可以给这个年轻的学院、给正在努力做事的人，也给我们自己多一点耐心。批评和期待并不矛盾，看到不足与相信事情能够慢慢变好，同样可以同时存在。}

\SourceParagraph{事实上，这本手册本身也是这种想法的一部分。}

\SourceParagraph{我们没有能力一下子改变学院，也没有资格替后来的人决定应该怎样度过大学生活。但至少，我们可以把自己已经知道的东西写下来，把一些有用的信息传递下去，把曾经得到过的帮助继续交给后来的人。如果每一届学生都愿意在前人的基础上再补上一点、修正一点、分享一点，那么几年以后，这份手册或许会成为学院成长过程中一份很有意思的记录。}

\SourceParagraph{所以，如果你在阅读这本手册时发现错误，欢迎指出；如果你有不同的经验，欢迎补充；如果你踩过我们没有踩过的坑，也欢迎写下来提醒后来的人。我们尤其希望，有一天你不仅是这本手册的读者，也会成为它的作者。}

\SourceParagraph{学院还很年轻，这本手册也一样。}

\SourceParagraph{希望它能够陪你更好地“生存”，但更希望若干年后回头看时，你会发现自己在这里做的事情，远不只是生存而已。}
\begin{flushright}
\begin{tabular}{@{}c@{}}
2026 年 8 月 20 日\\[1em]
编者按
\end{tabular}
\end{flushright}
\newpage

\addcontentsline{toc}{chapter}{免责声明}
\chapter*{免责声明}
\SourceParagraph{欢迎阅读《难喝生存手册 》！}
\SourceParagraph{本手册是若干 23、24、25 级南赫学生自发组织、共同撰写编写的非官方参考资料，其版权归属本指南编写小组。在您阅读、使用或分享本指南的任何内容前，请务必仔细阅读并充分理解以下条款。您继续使用本指南的行为，即表示您已同意并接受本声明的全部内容。}
\SourceListItem{1，适用范围与目的：手册旨在提供一般性帮助与指导，帮助读者了解相关信息与思路；本手册不应被解读为官方立场、承诺或专业意见。}
\SourceListItem{2，不构成专业意见：本手册的内容不构成任何专业意见、建议或服务，包括但不限于法律、投资、职业咨询等。如需相关专业领域的意见，请咨询具备资质的专业人士或机构。}
\SourceListItem{3，信息的准确性与时效性：我们尽力确保信息的准确性、完整性与时效性，但无法保证信息始终完全正确、最新或适用于所有情境。对因依赖本指南信息而产生的直接或间接损失不承担任何责任。}
\SourceListItem{4，使用风险与自我判断：使用本手册信息时，请结合自身情况自行判断并决定是否采取相关行动。因依赖本手册信息导致的决策及其后果，由读者自行承担全部风险与责任。}
\SourceListItem{5，第三方链接与资源：本手册可能包含指向第三方网站或资源的链接。对这些链接的内容、隐私政策及可用性不承担任何责任，访问与使用风险自行承担。}
\SourceListItem{6，知识产权与转载：本指南的文本、图片、图表等未获授权不得以商业用途转载、传播、再发行或衍生使用；如需转载，请标注出处并遵循相应的授权与引用规范。}
\SourceListItem{7，使用约束与行为规范：使用本手册应遵守相关法律法规、校园规章及学术诚信要求。不得利用本手册从事违法、侵权、诽谤、骚扰或其他不当行为。}
\SourceListItem{8，变更与更新：本手册保留随时修改、更新、删除或暂停发布的权利，恕不另行通知。以最新发布版本为准。}
\SourceParagraph{使用本指南即视为你已阅读、理解并同意遵守本免责声明的全部内容。若未来对免责声明条款进行修改，编写方将通过正式渠道发布更新版本，继续使用即表示接受更新后的条款。}
\newpage

\addcontentsline{toc}{chapter}{常用信息}
\chapter*{常用信息（各类可能用到的公众号、论坛、群聊等）}
\SourceParagraph{这些信息一个个细细翻完大概也要一周吧，希望大家能有所收获。}
\begin{center}
\small
\setlength{\tabcolsep}{2.5pt}
\renewcommand{\arraystretch}{1.22}
\begin{longtable}{|>{\RaggedRight\arraybackslash}p{0.08\textwidth}|>{\RaggedRight\arraybackslash}p{0.16\textwidth}|>{\RaggedRight\arraybackslash}p{0.10\textwidth}|>{\RaggedRight\arraybackslash}p{0.33\textwidth}|>{\RaggedRight\arraybackslash}p{0.20\textwidth}|}
\hline
生存手册相关 & 南哪助手 & 微信公众号 & 862670709（QQ群） & 语雀文档必看 \\
\hline
 & 上海交通大学生存手册 & 网页 & \url{https://survivesjtu.gitbook.io/survivesjtumanual/} & 经典无出其右 \\
\hline
 & 复旦生存手册 & 网页 & \url{https://fudanmanual.github.io/FudanManual/Intro/} & Empty较多；食之无味，弃之可惜 \\
\hline
 & 西安交大生存指南 & 网页 & \url{https://djm-xjtu.github.io/XJTU-Survival-Manual/} & 润学思想大纲格外有趣 \\
\hline
 & 武汉科技大学生存手册 & 网页 & \url{https://www.survivewustmanual.cn/} & 序很不错，少年心气 \\
\hline
 & 技科生存指南 & 网页 & \url{https://survive-jk-suzhoucompus.gitbook.io/survive-jk-suzhoucampus} & 从鼓楼校区到苏州校区 \\
\hline
大气专业相关 & 南赫资料分享（民间大群） & QQ群 & 1040472048 & 加入我们 \\
\hline
 & 气象家园 & 论坛 &  & 升学求职板块值得一看 \\
\hline
 & 气象学家 & 微信公众号 &  & 气象资讯一览无遗 \\
\hline
 & 气象人才 & 微信公众号 &  & 官方招聘信息发布平台 \\
\hline
 & EarthAi & 微信公众号 &  & 人工智能气象一线吃瓜 \\
\hline
 & 16度的可乐 & 小红书 & 256474813 & 气象就业n弹系列可以看看 \\
\hline
出国相关/交流交换项目 & 南哪儿留学 & 微信公众号 &  & 交换信息早知道 \\
\hline
 & 南京大学本科生院交换生管理系统 & 网页 & \url{http://elite.nju.edu.cn/exchangesystem/} & 交流项目官方平台 \\
\hline
 & NJU留学交流群 & QQ群 & 981368919（二群） & 不懂就问 \\
\hline
 & 南方科技大学飞跃手册 & 网页 & \url{https://sustech-application.com/} & 各类经验分享（神人云集，国内一流） \\
\hline
 & 清华大学飞跃手册 & 网页 & \url{https://feiyue.online/} & 不管是否想出国都非常推荐阅读；如切如磋，如琢如磨 \\
\hline
 & 南京大学各院系飞跃手册 & Pdf & 自己找 & 大气在历史上没有飞跃手册，我们申请结束后会做的。其他院系的手册也值得一看 \\
\hline
 & 一亩三分地社区 & 网页 & \url{https://www.1point3acres.com/} & 入门必看 经典无出其右 \\
\hline
大气出国 & 南大就业 & 微信公众号 & 2017 - 2024 南京大学大气科学学院就业报告（升学定位） & 大气出国必看 \\
\hline
 & 一展云图 & 微信公众号 &  & 国内外大气科学、气象学就业升学信息 \\
\hline
 & NJU大气留学交流群 & 微信群 &  & 讨论答疑（进群会被拷问） \\
\hline
 & 大气出国（2018年及以后） & QQ群 & 536999212 & （NJU校外）华人AP招生 \\
\hline
保研 & 南京大学南赫学院 & 微信公众号 &  & 你值得拥有 \\
\hline
 & 南京大学大气科学学院 & 微信公众号 &  & 招生、就业信息推送 \\
\hline
 & 清华大学地球系统科学系 & 微信公众号 &  & 招生、论坛信息推送 \\
\hline
 & 北京大学大气与海洋科学系 & 微信公众号 &  & 招生、暑校信息推送 \\
\hline
 & 中国科学院大气所研究生会 & 微信公众号 &  & 招生、就业信息推送 \\
\hline
 & 浙大地科学院 & 微信公众号 &  & 招生信息推送 \\
\hline
 & 和风沐雨光华复旦 & 微信公众号 &  & 招生、暑校信息推送 \\
\hline
 & 中国科学技术大学大气科学专业 & 微信公众号 &  & 招生、论坛信息推送 \\
\hline
 & 南京信息工程大学大气科学学院 & 微信公众号 &  & 招生、就业信息推送 \\
\hline
公考选调 & 等待国家分配工作的NJUers & QQ群 & 161075315 & NJU公考选调交流 \\
\hline
 & 南京大学基层研究会 & 微信公众号 &  & 选调相关信息分享 \\
\hline
\end{longtable}
\end{center}